\documentclass[11pt,a4paper]{article}
\usepackage[margin=2.2cm]{geometry}
\usepackage{amsmath,amssymb}
\usepackage[T1]{fontenc}
\usepackage{lmodern}
\usepackage{microtype}
\usepackage{xcolor}
\usepackage{enumitem}
\usepackage[colorlinks=true,urlcolor=blue!50!black]{hyperref}
\setlength{\parindent}{0pt}
\setlength{\parskip}{5pt}
\definecolor{qbar}{RGB}{90,110,160}
\newenvironment{comment}[1]{\par\medskip\noindent\textbf{Comment #1.}\par\vspace{2pt}%
  \begin{list}{}{\leftmargin=12pt\rightmargin=12pt}\item\itshape\small}{\end{list}\vspace{-2pt}}
\newcommand{\reply}{\noindent\textbf{Reply.} }
\newcommand{\changes}{\par\noindent\textbf{Changes.} }
\newcommand{\mf}[1]{\mathbb{#1}}
\begin{document}

\begin{center}
{\large\bfseries Reply to the referee report on JHEP\_282P\_0726}\\[3pt]
\emph{Soft gluon wave function and evolution operator in the CGC at next-to-leading order}\\[2pt]
Ramkumar Radhakrishnan
\end{center}
\vspace{-4pt}\hrule\vspace{6pt}

Dear Editor,

I thank the referee for a careful reading of the manuscript and for comments that have substantially improved it. Below I reply to each point. All changes are highlighted in colour in the revised manuscript. Equation numbers refer to the revised version unless marked ``v1''; in particular Eq.~(4.78) of v1 is now Eq.~(4.82), and Eq.~(4.82) of v1 is now Eq.~(4.88).

%------------------------------------------------------------------
\begin{comment}{1}
It seems that in Sec 4.3.1 one is neglecting a ``t-channel gluon exchange'' contribution, where one passes through a three-gluon intermediate state. This is a significant part of the result. It is also curious that the authors do not notice the absence of this contribution in the unitarity relations, which should not hold when it is neglected.
\end{comment}

\reply The referee is right. Eq.~(4.88) contains, in each of the three colour structures, the four-gluon contact interaction and the \emph{instantaneous} gluon exchange in the $t$-, $u$- and $s$-channel, which comes from the instantaneous term of the light-cone Hamiltonian. The \emph{propagating} exchange, which is second order in $H_{ggg}$ and passes through an intermediate state, was included only in the $s$-channel, Eq.~(4.90). The propagating $t$- and $u$-channel exchanges, with a three-gluon intermediate state, were missing.

They are now computed explicitly in a new part of Sec.~4.3.1, ``$t$- and $u$-channel gluon exchange'', Eqs.~(4.91)--(4.94) and Fig.~2. Each channel has two light-cone time orderings, depending on which incoming gluon emits the exchanged gluon. For given external momenta exactly one of them contributes, which is made explicit by $\Theta$-functions. The intermediate energy denominators vanish only for collinear splitting, where the vertex factors vanish linearly, so the transverse integrals are finite. The four corresponding terms are added to the coefficient $\mf{D}_{2}$ in Eq.~(5.7).

The referee is also right about unitarity. For the intermediate state $|g(r)g(r')g(p)\rangle$ the unitarity identity, Eq.~(4.15), produces products $-\mf{C}^{\dagger}(q,r',p)\,\mf{C}(k,r',r)$, in which the internal momentum connects an incoming and an outgoing gluon on different lines. In v1 these products had no counterpart in $\mf{D}_{2}$; they are reproduced exactly by the new terms. The revised text states that $\mf{B}_{3}$ and $\mf{D}_{2}$ are fixed completely by the matching to the light-cone wave function, so that their unitarity relations are non-trivial checks of the calculation, and it shows how these checks are carried out, Eq.~(4.15).

While redoing this part I also corrected the $s$-channel result, Eq.~(4.90): a factor $1/r'^{+}$ and factors of 2 in the intermediate energy denominator were missing, and the colour factor and one Kronecker delta are now consistent with the vertex, Eq.~(D.3).

\changes Paragraph after Eq.~(4.88); corrected Eq.~(4.90); new Eqs.~(4.91)--(4.94) and Fig.~2; four new terms in Eq.~(5.7); new paragraph with Eq.~(4.15) and remark after Eq.~(5.12).

%------------------------------------------------------------------
\begin{comment}{2}
This is related to both the ``t-channel gluon exchange'' contribution mentioned above and the contribution already included in (4.82). The $2g\to2g$ energy denominator can go to zero and thus needs to be regularized. How this is done should be discussed in the manuscript.
\end{comment}

\reply I agree that this must be stated. Even after the momentum-conserving $\delta$-function is used, the denominator $E_{i}-E_{f}$ of Eqs.~(4.88) and (4.91)--(4.94) vanishes on a set of measure zero, e.g.\ for forward kinematics $r=p$, $q=k$.

The wave functions are computed in non-degenerate (Rayleigh--Schr\"odinger) perturbation theory. Since the free soft-gluon states form a continuum, the degenerate configurations have measure zero, and the series is defined in the standard way by adiabatic switching. In the revised manuscript $\Omega$ is defined as the interaction-picture evolution operator $U_{\epsilon}(0,-\infty)$, Eq.~(3.7), i.e.\ the M{\o}ller operator of scattering theory. This gives the prescription
\begin{equation*}
E_{i}-E_{f}\;\to\;E_{i}-E_{f}+in\epsilon \quad\text{at $n$-th order},\qquad \epsilon\to0^{+},
\end{equation*}
Eq.~(3.8). With this prescription the unitarity relations, Eq.~(3.9), hold identically for every $\epsilon>0$. A pole at $E_{i}=E_{f}$ remains only in $\mf{B}_{3}$ and in the $\rho$-independent part of $\mf{D}_{2}$, where its residue is the on-shell $2\to2$ scattering amplitude. Using $1/(x+i\epsilon)=\mathrm{P}(1/x)-i\pi\delta(x)$, the principal value is integrable because the pole is simple, and the $\delta$-function term is an on-shell, anti-Hermitian contribution that drops out of the unitarity relations.

The diagonalization of Sec.~6 is unaffected. In $\Omega^{\dagger}H\Omega$ the off-diagonal remainder is multiplied by $in\epsilon/(E_{i}-E_{f}+in\epsilon)$. This factor vanishes as $\epsilon\to0$ everywhere except on the energy shell, which has measure zero, so it does not contribute to matrix elements between normalizable states.

\changes New paragraphs in Sec.~3.2; Eq.~(3.7) rewritten; new Eqs.~(3.8)--(3.9); a sentence after Eq.~(4.88); Appendix~E rewritten in the standard Rayleigh--Schr\"odinger form.

%------------------------------------------------------------------
\begin{comment}{3}
At this order in perturbation theory the QCD coupling should run and the NLO result should depend on a scale associated with the coupling constant renormalization. Indeed one has here a $\beta$-function coefficient in (4.78), but somehow no running of the coupling constant. This should be addressed in the text.
\end{comment}

\reply The constant $11/6$ in Eq.~(4.82) multiplies the transverse integral $\int d^{2}\tilde p/\tilde p^{2}$, which is scaleless and vanishes in dimensional regularization. It vanishes because its UV and IR divergences cancel, not because it is finite:
\begin{equation*}
\mu^{2\epsilon}\!\int\!\frac{d^{2-2\epsilon}\tilde p}{\tilde p^{2}}=\pi\Big(\frac{1}{\epsilon_{\rm UV}}-\frac{1}{\epsilon_{\rm IR}}\Big)=0 .
\end{equation*}
The UV part of the $11/6$ term is $-\dfrac{g^{2}}{16\pi^{2}}\dfrac{\beta_{0}}{\epsilon_{\rm UV}}$ with $\beta_{0}=\tfrac{11}{3}N_{c}$. This is the one-loop gluon field-strength renormalization, which in light-cone gauge also fixes the coupling renormalization, $Z_{g}=Z_{3}^{-1/2}$ [Leibbrandt, Rev.\ Mod.\ Phys.\ 59 (1987) 1067]. This explains why the $\beta$-function coefficient appears.

Hence the normalization $\mf{N}$ contains no renormalization scale at $\mathcal{O}(g^{2})$; the UV divergence is absorbed into the renormalized coupling $g(\mu)$. The scale enters through the coupling. The $\mathcal{O}(g)$ coefficients $\mf{A}$ and $\mf{C}$ are tree level and are functions of $g(\mu)$. Their one-loop, $\mathcal{O}(g^{3})$, corrections [Lublinsky \& Mulian, JHEP 05 (2017) 097] carry the compensating $\log\mu^{2}$ with coefficient $\propto\beta_{0}$. Together these turn the fixed coupling of the Weizs\"acker--Williams amplitude into a running coupling at a scale set by the gluon transverse momentum, as in the running-coupling BK/JIMWLK kernels [Balitsky, PRD 75 (2007) 014001; Kovchegov \& Weigert, NPA 784 (2007) 188]. Computing these $\mathcal{O}(g^{3})$ terms is beyond the $\mathcal{O}(g^{2})$ evolution operator constructed here.

\changes Text after Eq.~(4.82) and new Eqs.~(4.83)--(4.85); a sentence after Eq.~(5.2); the corresponding remark after Eq.~(6.8) and in the Conclusions; three new references.

%------------------------------------------------------------------
\begin{comment}{4}
In Sec 3.1 it is stated that the ``coefficients are not ordinary c-number functions but operator quantities acting in the valence Hilbert space''. [\ldots] The statement about noncommuting operator coefficients only applies to the soft gluon part of the Hilbert space, not the full one. The fact that the states in e.g.\ eq (3.2) are states in only the soft subspace should be emphasized more in the text.
\end{comment}

\reply I agree, and the text has been rewritten. The revised Sec.~3.1 states that Eq.~(3.2) is an expansion in the Fock basis of the soft subspace $\mathcal{H}_{\rm soft}$ only. The full Hilbert space $\mathcal{H}_{\rm soft}\otimes\mathcal{H}_{\rm val}$ is an ordinary complex vector space, and in a complete basis $|n\rangle_{\rm soft}\otimes|v\rangle_{\rm val}$ all coefficients are $c$-numbers. Because the valence part is not expanded, the coefficients of the soft Fock states are functionals of the valence charge density, i.e.\ operators on $\mathcal{H}_{\rm val}$, and their ordering matters because of Eq.~(2.19). All states of the form (3.2), and the vacuum $|0\rangle\equiv|0\rangle_{\rm soft}\otimes|v\rangle$, are to be understood in this sense. The same remark is repeated where $\Omega$ is expanded in Sec.~4.

\changes Rewritten paragraph after Eq.~(3.2); sentence after Eq.~(4.1).

%------------------------------------------------------------------
\medskip
\textbf{Other changes.} In a complete re-check of the manuscript I also corrected typographical errors in a number of equations: signs, indices, integration measures, and missing Hermitian-conjugate terms in Secs.~2, 4, 5 and 6 and in Appendices~A--C. I also stated the plus-momentum window that regularizes the gluon-number-preserving transitions in Sec.~4.2. In addition, I clarified that the $\mathcal{O}(g)$ components of $\Omega|gg\rangle$ and $\Omega|ggg\rangle$ contribute to the normalization and that $\mf{N}$ is fixed by unitarity, and stated the terms that cancel in the diagonalization, after Eq.~(6.7). The Introduction now refers to pure Yang--Mills theory rather than QCD, consistent with the content of the paper.

I hope that the revised manuscript is now suitable for publication in JHEP.

\medskip
Sincerely,\\
Ramkumar Radhakrishnan
\end{document}
